% Source: https://artofproblemsolving.com/community/c3h2067835_math_help_pls
% Stable identity: aops-topic-2067835-q01
(a) ``$x$ varies directly as $y$ and inversely as $z^2$'' means there is one constant $K$ with
\[
x=\dfrac{Ky}{z^2},\qquad\text{equivalently}\qquad \dfrac{xz^2}{y}=K\ \text{ is constant.}
\]
From the data, $K=\dfrac{10\cdot196}{4}=490$. For $y=16$ and $z=7$: $x=\dfrac{490\cdot16}{49}=\boxed{160}$.

(b) Joint variation does \emph{not} mean that $x$ is proportional to $y$ for all values of $z$ and, separately, inversely proportional to $z^2$ for all values of $y$. ``$x$ varies directly as $y$'' holds only while $z$ is fixed, and ``inversely as $z^2$'' only while $y$ is fixed. The student's relation $kx=y$ ignores $z$, which changes from $14$ to $7$, and $\dfrac {c}{x}=z^2$ ignores $y$, which changes from $4$ to $16$. The two separate equations are not consistent with the original statement (both cannot hold with fixed constants when both $y$ and $z$ change). The correct way is to combine the two effects: multiplying $y$ by $4$ multiplies $x$ by $4$, and halving $z$ multiplies $x$ by $2^2=4$, so $x=10\cdot4\cdot4=160$.
